\documentclass[11pt]{article}
\usepackage[margin=1in]{geometry}
\usepackage{amsmath,amssymb,amsthm,hyperref}
\newtheorem{theorem}{Theorem}
\newtheorem{lemma}{Lemma}
\newcommand{\fl}[1]{\left\lfloor#1\right\rfloor}
\newcommand{\fr}[1]{\left\{#1\right\}}
\title{The four limiting values of a square-root floor sum}
\author{Principia Math In Progress}
\date{29 September 2026}
\begin{document}\maketitle
\begin{abstract}
We prove all four residue-class limits conjectured by Chamberland and
Dilcher. An elementary complete quadratic exponential-sum bound gives
equidistribution of the normalized square remainders for every sequence
of positive integer moduli. An exact lattice count then transfers their
mean to the floor sum. No restriction on prime factors or square factors
is required.
\end{abstract}

\section*{Task 1: the conjecture and its complete resolution}
Following~\cite[Equation (1.4)]{CD}, define
\[
 f(n)=\sum_{j=1}^{\lfloor n/4\rfloor}\fl{\sqrt{jn}}-\frac{n^2-1}{12}
 \qquad(n\ge1).
\]
\begin{theorem}\label{thm:main}
For each $r\in\{0,1,2,3\}$,
\[
 \lim_{m\to\infty}\frac{f(4m+r)}{4m+r}=\frac{1-r}{8}.
\]
Thus the four limits, in order, are $1/8$, $0$, $-1/8$, and $-1/4$.
\end{theorem}

\paragraph{The last numerator in the source.}
Conjecture 7.5, Equation (7.14), prints $f(4m+2)/(4m+3)$ with
limit $-1/4$. Its intended numerator is $f(4m+3)$: the discussion
immediately below cites identities (1.6), (6.5), and (7.6), all concerning
the residue class $3$ modulo $4$.
The printed expression instead has limit $-1/8$, by the third case of
Theorem~\ref{thm:main} and $(4m+2)/(4m+3)\to1$.
We prove the full intended four-case assertion, as well as resolve the
literal printed expression; no credit is claimed merely for this typo.

\section*{Mean of the normalized square remainders}
Write $e(t)=\exp(2\pi i t)$ and
\[
 G_n(h)=\sum_{k=0}^{n-1}e(hk^2/n),\qquad h\in\mathbb Z\setminus\{0\}.
\]
Changing variables $k=\ell+u$ modulo $n$ gives
\[
 |G_n(h)|^2
 =\sum_{u\bmod n}e(hu^2/n)\sum_{\ell\bmod n}e(2hu\ell/n).
\]
The inner sum is $n$ when $n\mid2hu$ and zero otherwise. There are
$\gcd(2h,n)$ such residues $u$, so
\begin{equation}\label{eq:gauss}
 |G_n(h)|^2\le n\gcd(2h,n)\le2|h|n.
\end{equation}
Consequently $G_n(h)/n\to0$ as $n\to\infty$ for each fixed nonzero $h$.
By linearity, for every trigonometric polynomial $P$,
\[
 \frac1n\sum_{k=0}^{n-1}P(k^2/n)\longrightarrow\int_0^1P(t)\,dt.
\]
The same assertion holds for every continuous periodic function $g$,
by uniform approximation with trigonometric polynomials.
For completeness, this approximation follows by convolving $g$ with
the Fej\'er kernels
\[
 K_M(t)=\frac1M\left|\sum_{j=0}^{M-1}e(jt)\right|^2.
\]
They are nonnegative trigonometric polynomials of integral one, and
their mass outside any fixed neighborhood of the integers tends to zero:
if the distance to the nearest integer is at least $\eta\in(0,1/2)$,
then $K_M(t)\le1/(M\sin^2(\pi\eta))$.
Uniform continuity of $g$ therefore gives uniform convergence of
$K_M*g$ to $g$.

We must also handle the discontinuity of the fractional-part function.
Fix $0<\varepsilon<1/2$. A continuous periodic majorant $U_\varepsilon$
of $t$ on $[0,1)$ is obtained by using $U_\varepsilon(t)=t$ for
$\varepsilon\le t<1$ and joining $(0,1)$ to $(\varepsilon,\varepsilon)$
linearly. A continuous periodic minorant $L_\varepsilon$ equals $t$ on
$[0,1-\varepsilon]$ and joins $(1-\varepsilon,1-\varepsilon)$ to $(1,0)$
linearly. Their integrals are $1/2+\varepsilon/2$ and
$1/2-\varepsilon/2$, respectively. Applying the continuous-function
limit to both envelopes and then letting $\varepsilon\downarrow0$ proves
\begin{equation}\label{eq:mean}
 \frac{S_n}{n}\longrightarrow\frac12,\qquad
 S_n:=\sum_{k=0}^{n-1}\fr{\frac{k^2}{n}}
     =\frac1n\sum_{k=0}^{n-1}(k^2\bmod n).
\end{equation}
This includes all positive moduli and also includes the residues whose
squares are divisible by $n$.

\section*{An exact lattice identity with square endpoints}
Put $r=n\bmod4$, $\nu=\lfloor n/4\rfloor$, $m=2\nu=(n-r)/2$,
and $R_n(k)=k^2\bmod n$. Since
$\lfloor\sqrt{\nu n}\rfloor=2\nu$, counting pairs with
$1\le j\le\nu$ and $k^2\le jn$ gives
\begin{align*}
 \sum_{j=1}^{\nu}\fl{\sqrt{jn}}
 &=\sum_{k=1}^{m}\left(\nu+1-\left\lceil\frac{k^2}{n}\right\rceil\right)\\
 &=2\nu^2-\sum_{k=1}^{m}\fl{\frac{k^2}{n}}+Z_n,
\end{align*}
where $Z_n=\#\{k:1\le k\le m,\ n\mid k^2\}$.
The identity remains valid for $\nu=0$, with empty sums.
Define
\[
 A_n=\frac1n\sum_{k=1}^{m}R_n(k).
\]
Replacing each floor by the square minus its remainder, summing the
consecutive squares, and subtracting $(n^2-1)/12$ yields the exact formula
\begin{equation}\label{eq:exact}
 f(n)=A_n-\frac{r+1}{8}n+\frac r4+
              \frac{r(r-1)(r-2)}{24n}+Z_n.
\end{equation}

The errors in replacing $A_n$ by $S_n/2$ and ignoring $Z_n$ are
sublinear. First, $R_n(n-k)=R_n(k)$. The range $1\le k\le m$
therefore gives half the full remainder sum, except for at most one
central term or half of that term. Thus
\begin{equation}\label{eq:half}
 |A_n-S_n/2|\le1.
\end{equation}
For example, when $n$ is even, $m$ is either $n/2$ or $n/2-1$,
and the difference is $\pm R_n(n/2)/(2n)$. When $n$ is odd, $m$
is $(n-1)/2$ or $(n-3)/2$, giving difference zero or
$-R_n((n-1)/2)/n$.

Next write $n=\prod_p p^{a_p}$ and put
$D_n=\prod_p p^{\lceil a_p/2\rceil}$. Then $n\mid k^2$ precisely
when $D_n\mid k$, so
\[
 0\le Z_n\le\frac n{D_n}
 =\prod_p p^{\lfloor a_p/2\rfloor}\le\sqrt n.
\]
Since $0\le r/4+r(r-1)(r-2)/(24n)\le1$, equations
\eqref{eq:exact} and~\eqref{eq:half} imply, for every positive integer $n$,
\[
 \left|f(n)-\frac{S_n}{2}+\frac{r+1}{8}n\right|
 \le\sqrt n+2.
\]
Divide by $n$ and use~\eqref{eq:mean}. For each fixed residue $r$ this gives
\[
 \frac{f(n)}n\longrightarrow\frac14-\frac{r+1}{8}
 =\frac{1-r}{8}\qquad(n\to\infty,\ n\equiv r\pmod4),
\]
proving Theorem~\ref{thm:main}.

\section*{Tasks 2--5: follow-up questions}
\begin{enumerate}\setcounter{enumi}{1}
\item Determine the sharp uniform rate of convergence of $S_n/n$ to
$1/2$, including the dependence on the prime factorization of $n$.
\item Determine the order and limiting behavior of
$f(n)-(1-r)n/8$ on each residue class, after the best natural rescaling.
\item Find infinite families of moduli on which the square-endpoint term
$Z_n$ or the quadratic-remainder discrepancy determines the main error.
\item Extend the argument to floor sums of higher roots and identify
the residue-class main terms and the required exponential-sum bounds.
\end{enumerate}

\paragraph{Verification and scope.}
The verifier checks the exact lattice identity, square-endpoint count,
and comparison with the original floor sum using integer and rational
arithmetic. The limiting argument is proved analytically above; numerical
examples do not establish convergence. All prime factorizations are
allowed. Review is by the author-agent, not independent peer review.

\begin{thebibliography}{9}
\bibitem{CD} M.~Chamberland and K.~Dilcher,
\emph{Sums of the floor function related to class numbers of imaginary
quadratic fields},
\href{https://arxiv.org/abs/2510.04387}{arXiv:2510.04387v1} (2025),
Conjecture 7.5, p.~23. All 23 pages of the current manuscript were read
on 29 September 2026.
\bibitem{MathDB} MathDB,
\href{https://mathdb.com/p/370789}{problem 370789},
live status checked 29 September 2026.
\end{thebibliography}
\end{document}
